% GENERATED FILE. Edit canonical lessons, then run npm run generate:books.
% canonical-source-hash: fnv1a64:784806839c115c62
% canonical-lessons: BN-C191-dorkari, BN-C191-sombhob, BN-C191-adhunik, BN-C191-prachin, BN-C191-asol

\chapter{Necessary, Possible, Modern, Ancient, Real}
\label{ch:bn-necessary-possible-modern-ancient-real}
\hlchaptermodality{191}

\begin{chapteropening}
\textbf{By the end of this chapter:} \emph{I can say necessary, possible, modern, ancient and real.}

Complete the last lesson of chapter 191: I can say necessary, possible, modern, ancient and real.
\end{chapteropening}

\section[\texorpdfstring{\bn{দরকারি} (dôrkāri)}{দরকারি (dôrkāri)}]{\bn{দরকারি} (dôrkāri) — necessary}

\label{lesson:BN-C191-dorkari}



\begin{quote}
Before the new one: say the Bengali for unfortunately, then the Bengali for a result.
\end{quote}

\subsection*{You'll want to know: \bn{দরকারি}}
\textbf{\bn{দরকারি}} — \emph{dôrkāri} — \textquotedblleft{}necessary\textquotedblright{}.

\begin{grammarlens}[title={What you've built}]
One more word for this chapter.
\end{grammarlens}

\subsection*{Your turn}
\begin{itemize}
\raggedright
  \item \emph{Say it:} \emph{dôrkāri}
  \item \emph{Say it:} \emph{dôrkāri}, once more
  \item \emph{Say it:} say \emph{dôrkāri}
  \item \emph{Recall:} say the Bengali for unfortunately, then the Bengali for a result, then say \emph{dôrkāri} again
\end{itemize}

\begin{tcolorbox}[breakable,colback=teal!4,colframe=teal!35!black,title={Before you move on}]
What does \bn{দরকারি} mean? (\textquotedblleft{}Necessary\textquotedblright{}.) Say it once more.
\end{tcolorbox}
\section[\texorpdfstring{\bn{সম্ভব} (sômbhôb)}{সম্ভব (sômbhôb)}]{\bn{সম্ভব} (sômbhôb) — possible}

\label{lesson:BN-C191-sombhob}



\begin{quote}
Before the new one: say the Bengali for a result, then the Bengali for necessary.
\end{quote}

\subsection*{You'll want to know: \bn{সম্ভব}}
\textbf{\bn{সম্ভব}} — \emph{sômbhôb} — \textquotedblleft{}possible\textquotedblright{}.

\begin{grammarlens}[title={What you've built}]
One more word for this chapter.
\end{grammarlens}

\subsection*{Your turn}
\begin{itemize}
\raggedright
  \item \emph{Say it:} \emph{sômbhôb}
  \item \emph{Say it:} \emph{sômbhôb}, once more
  \item \emph{Say it:} say \emph{sômbhôb}
  \item \emph{Recall:} say the Bengali for a result, then the Bengali for necessary, then say \emph{sômbhôb} again
\end{itemize}

\begin{tcolorbox}[breakable,colback=teal!4,colframe=teal!35!black,title={Before you move on}]
What does \bn{সম্ভব} mean? (\textquotedblleft{}Possible\textquotedblright{}.) Say it once more.
\end{tcolorbox}
\section[\texorpdfstring{\bn{আধুনিক} (ādhunik)}{আধুনিক (ādhunik)}]{\bn{আধুনিক} (ādhunik) — modern}

\label{lesson:BN-C191-adhunik}



\begin{quote}
Before the new one: say the Bengali for necessary, then the Bengali for possible.
\end{quote}

\subsection*{You'll want to know: \bn{আধুনিক}}
\textbf{\bn{আধুনিক}} — \emph{ādhunik} — \textquotedblleft{}modern\textquotedblright{}.

\begin{grammarlens}[title={What you've built}]
One more word for this chapter.
\end{grammarlens}

\subsection*{Your turn}
\begin{itemize}
\raggedright
  \item \emph{Say it:} \emph{ādhunik}
  \item \emph{Say it:} \emph{ādhunik}, once more
  \item \emph{Say it:} say \emph{ādhunik}
  \item \emph{Recall:} say the Bengali for necessary, then the Bengali for possible, then say \emph{ādhunik} again
\end{itemize}

\begin{tcolorbox}[breakable,colback=teal!4,colframe=teal!35!black,title={Before you move on}]
What does \bn{আধুনিক} mean? (\textquotedblleft{}Modern\textquotedblright{}.) Say it once more.
\end{tcolorbox}
\section[\texorpdfstring{\bn{প্রাচীন} (prāchīn)}{প্রাচীন (prāchīn)}]{\bn{প্রাচীন} (prāchīn) — ancient}

\label{lesson:BN-C191-prachin}



\begin{quote}
Before the new one: say the Bengali for possible, then the Bengali for modern.
\end{quote}

\subsection*{You'll want to know: \bn{প্রাচীন}}
\textbf{\bn{প্রাচীন}} — \emph{prāchīn} — \textquotedblleft{}ancient\textquotedblright{}.

\begin{grammarlens}[title={What you've built}]
One more word for this chapter.
\end{grammarlens}

\subsection*{Your turn}
\begin{itemize}
\raggedright
  \item \emph{Say it:} \emph{prāchīn}
  \item \emph{Say it:} \emph{prāchīn}, once more
  \item \emph{Say it:} say \emph{prāchīn}
  \item \emph{Recall:} say the Bengali for possible, then the Bengali for modern, then say \emph{prāchīn} again
\end{itemize}

\begin{tcolorbox}[breakable,colback=teal!4,colframe=teal!35!black,title={Before you move on}]
What does \bn{প্রাচীন} mean? (\textquotedblleft{}Ancient\textquotedblright{}.) Say it once more.
\end{tcolorbox}
\section[\texorpdfstring{\bn{আসল} (āsôl)}{আসল (āsôl)}]{\bn{আসল} (āsôl) — real, genuine}

\label{lesson:BN-C191-asol}



\begin{quote}
Before the new one: say the Bengali for modern, then the Bengali for ancient.
\end{quote}

\subsection*{You'll want to know: \bn{আসল}}
\textbf{\bn{আসল}} — \emph{āsôl} — \textquotedblleft{}real, genuine\textquotedblright{}.

\begin{grammarlens}[title={What you've built}]
One more word for this chapter.
\end{grammarlens}

\subsection*{Your turn}
\begin{itemize}
\raggedright
  \item \emph{Say it:} \emph{āsôl}
  \item \emph{Say it:} \emph{āsôl}, once more
  \item \emph{Say it:} say \emph{āsôl}
  \item \emph{Recall:} say the Bengali for modern, then the Bengali for ancient, then say \emph{āsôl} again
\end{itemize}

\begin{tcolorbox}[breakable,colback=teal!4,colframe=teal!35!black,title={Before you move on}]
What does \bn{আসল} mean? (\textquotedblleft{}Real, genuine\textquotedblright{}.) Say it once more.
\end{tcolorbox}
